\chapter{Driven Simulations and S-Parameters}
\label{ch:driven}

\chapterorient{We have built the anchor: electrostatics assembles one operator,
\emph{hoists} it out of the loop, and maps a family of independent solves over it
(\cref{ch:electrostatics}). This chapter takes that anchor and moves a single
line of code---the line that captures the operator into the solver---from
\emph{outside} the loop to \emph{inside} it. That one move changes the analysis
from electrostatics to the driven frequency sweep, the workhorse of RF and
microwave engineering. We feed a time-harmonic signal into the ports of a device,
sweep the frequency across a band, and measure how much of each wave is
transmitted and how much reflected: the scattering matrix $S(\omega)$. Along the
way the operator is no longer fixed---it is rebuilt at every frequency from a
small fixed basis---and the reduction is no longer a Gram but a port-projection.
By the end you will see the driven analysis for exactly what it is: the anchor
with the operator rebuilt \emph{inside} the sweep, capped by a projection
readout.}

\section{The physical question: what a network analyzer measures}

Point a vector network analyzer at a microwave device---a filter, a coupler, a
length of transmission line, a resonator on a chip---and it does something very
concrete. It connects to the device's \emph{ports}: the places where signal
enters and leaves, typically coaxial connectors or on-chip waveguide
cross-sections. It then injects a pure sinusoid at one port, at one frequency
$\omega$, waits for the transients to die away so the device is in
\emph{time-harmonic steady state}, and measures two things at every port: how
much of the injected wave came \emph{back} out of the port it went in
(reflection), and how much came out of each \emph{other} port (transmission). It
repeats this across a whole band of frequencies. The result is a table of complex
numbers $S_{ij}(\omega)$---the \emph{scattering parameters}---and it is the
universal currency of the field. A filter's passband, a coupler's split ratio, an
antenna's match, a chip interconnect's loss: all of them are read directly off
$S(\omega)$.

\begin{physicsbox}
A two-port device (\cref{fig:setup}) sits between a source and a load. At a single
frequency, the analyzer launches an \emph{incident} wave $a_1$ into port~1. Part
of it \emph{reflects} straight back out of port~1 (the wave $b_1$), and part
\emph{transmits} through the device and emerges at port~2 (the wave $b_2$). The
scattering parameters are the ratios of outgoing to incoming wave amplitudes,
\[
  b_1 = S_{11}\,a_1 + S_{12}\,a_2, \qquad
  b_2 = S_{21}\,a_1 + S_{22}\,a_2,
\]
so that $S_{ij}$ is ``the wave that comes out of port~$i$ per unit wave injected
into port~$j$, with all other ports terminated.'' The diagonal entries are
\emph{reflection} coefficients (how much bounces back); the off-diagonals are
\emph{transmission} coefficients (how much gets through). Everything in this
chapter exists to compute that matrix, frequency by frequency.
\end{physicsbox}

\begin{figure}[t]
\centering
\begin{tikzpicture}[font=\footnotesize, >=Stealth]
  % the device
  \node[stage, minimum width=30mm, minimum height=18mm, fill=simBgGray,
        draw=simGray] (dev) at (0,0) {2-port device\\\scriptsize (filter / line)};
  % port 1 (left) feed line
  \draw[very thick, simBlue] (-3.6,0.35) -- (dev.west |- 0,0.35);
  \draw[very thick, simBlue] (-3.6,-0.35) -- (dev.west |- 0,-0.35);
  \node[font=\scriptsize, text=simBlue] at (-3.0,0.85) {port 1};
  % port 2 (right) feed line
  \draw[very thick, simGreen] (dev.east |- 0,0.35) -- (3.6,0.35);
  \draw[very thick, simGreen] (dev.east |- 0,-0.35) -- (3.6,-0.35);
  \node[font=\scriptsize, text=simGreen] at (3.0,0.85) {port 2};
  % incident wave a1
  \draw[->, thick, simRust] (-3.2,1.25) -- (-1.9,1.25)
     node[midway, above, font=\scriptsize]{incident $a_1$};
  % reflected wave b1
  \draw[->, thick, simRust] (-1.9,-1.25) -- (-3.2,-1.25)
     node[midway, below, font=\scriptsize]{reflected $b_1=S_{11}a_1$};
  % transmitted wave b2
  \draw[->, thick, simGreen!70!black] (1.9,1.25) -- (3.2,1.25)
     node[midway, above, font=\scriptsize]{transmitted $b_2=S_{21}a_1$};
  % the field inside
  \node[font=\scriptsize\itshape, text=simInk] at (0,0.0) {};
  \node[font=\scriptsize\itshape, align=center] at (0,-1.5)
    {drive port 1, terminate port 2:\\ read $S_{11}$ (reflect), $S_{21}$ (transmit)};
\end{tikzpicture}
\caption[A two-port device and its scattering waves]{The physical setup. A
time-harmonic wave $a_1$ is injected into port~1 of a two-port device. Part
reflects back out of port~1 (the wave $b_1 = S_{11}a_1$); part transmits through
and emerges at port~2 ($b_2 = S_{21}a_1$). Driving port~1 and reading both ports
gives the first column of the scattering matrix; driving port~2 gives the second.
A simulation reproduces this measurement: it solves the time-harmonic field
inside the device with port~1 excited, then projects the solved field onto each
port's mode to recover $S_{11}$ and $S_{21}$.}
\label{fig:setup}
\end{figure}

The analyzer's loop---inject at a frequency, measure the scattered waves, step to
the next frequency---is exactly the loop the simulation runs. The simulator does
not have a physical sinusoid and a physical analyzer; it has the time-harmonic
\emph{field equation} and a way to read scattered waves off the solved field. But
the shape is identical: a sweep over frequency, an independent steady-state
problem at each one, and a projection readout. That shape is the operator-varying
map of \cref{ch:situating}, and the rest of this chapter is its anatomy.

\section{Assemble the basis once: the frequency-dependent operator}
\label{sec:assemble}

Recall the time-harmonic reduction of Maxwell from \cref{ch:reductions-pde}. A
field oscillating at a single angular frequency $\omega$ is written
$\vE(\vx,t) = \Re\{\vE(\vx)\,e^{i\omega t}\}$, and the substitution
$\partial_t \mapsto i\omega$, $\partial_t^2 \mapsto -\omega^2$ turns Maxwell's
curl--curl equation with sources into a single complex linear system for the
phasor field $\vE(\vx)$. Discretized on a Nédélec $\Hcurl$ finite-element space
(\cref{ch:functionspaces}), that system reads
\begin{equation}
  \mat{A}(\omega)\,\vE \;=\; \vb(\omega),
  \qquad
  \mat{A}(\omega) \;=\; \Kop \;+\; i\omega\,\Cop \;-\; \omega^2\Mop
    \;\bigl[\,+\,\mat{A}_2(\omega)\,\bigr],
  \label{eq:driven-system}
\end{equation}
where the three operators are the same characters from every chapter of this
book: $\Kop$ is the curl--curl \emph{stiffness} (built from the reluctivity
$\nu=1/\mu$), $\Cop$ is the \emph{damping} (built from the surface impedance /
conductivity---the term that makes the system lossy and complex), and $\Mop$ is
the \emph{mass} (built from the permittivity $\varepsilon$). The optional
$\mat{A}_2(\omega)$ collects any genuinely frequency-dependent contribution that
is not a clean monomial in $\omega$---a dispersive material model, or a
wave-port term whose modal admittance varies with frequency.

Here is the structural fact this whole chapter turns on, and it is worth stating
in two halves.

\emph{The basis is fixed.} The three operators $\Kop$, $\Cop$, $\Mop$ do not
depend on $\omega$. They are built from the mesh, the function space, and the
materials---all of which are frozen the moment the configuration is parsed. So we
assemble them exactly \emph{once}, before the sweep ever begins, and store them as
a small fixed record (\cref{ch:records}):
\[
  \texttt{fam} \;=\; \{\,\Kop,\ \Cop,\ \Mop,\ \mat{A}_2\,\}.
\]
This is precisely the ``assemble once'' that electrostatics did with its single
$\Kop$ (\cref{ch:electrostatics}); the driven analysis just has three operators in
the basis instead of one.

\emph{The system operator is not fixed.} The operator that actually appears on the
left-hand side of the solve, $\mat{A}(\omega)$, is a frequency-dependent
\emph{combination} of the fixed basis. At each frequency it is a different
operator: the weights $1$, $i\omega$, $-\omega^2$ multiply the same three stored
operators, but those weights change with $\omega$, so $\mat{A}(\omega_1)$ and
$\mat{A}(\omega_2)$ are genuinely different matrices. We never store
$\mat{A}(\omega)$ permanently; we \emph{form} it from the basis at each frequency
we visit, use it for one solve, and discard it.

\begin{figure}[t]
\centering
\begin{tikzpicture}[font=\footnotesize, >=Stealth]
  % fixed basis on the left
  \node[opbox, minimum width=9mm] (K) at (0,1.4) {$\Kop$};
  \node[opbox, minimum width=9mm] (C) at (0,0.0) {$\Cop$};
  \node[opbox, minimum width=9mm] (M) at (0,-1.4) {$\Mop$};
  \node[font=\scriptsize\itshape, text=simGray, align=center] at (0,-2.3)
    {fixed basis\\(assembled once)};
  % weights
  \node[font=\footnotesize] (wK) at (2.7,1.4) {$1$};
  \node[font=\footnotesize] (wC) at (2.7,0.0) {$i\omega$};
  \node[font=\footnotesize] (wM) at (2.7,-1.4) {$-\omega^2$};
  \node[font=\scriptsize\itshape, text=simRust, align=center] at (2.7,-2.3)
    {weights\\(depend on $\omega$)};
  \draw[->, simGray] (K.east) -- (wK.west);
  \draw[->, simGray] (C.east) -- (wC.west);
  \draw[->, simGray] (M.east) -- (wM.west);
  % sum node
  \node[draw=simBlue, circle, thick, fill=simBgBlue, inner sep=1.5pt]
    (sum) at (5.0,0.0) {$\Sigma$};
  \draw[->, simInk] (wK.east) .. controls +(0.8,0) and +(-0.8,0.3) .. (sum.west);
  \draw[->, simInk] (wC.east) -- (sum.west);
  \draw[->, simInk] (wM.east) .. controls +(0.8,0) and +(-0.8,-0.3) .. (sum.west);
  % result
  \node[product, minimum width=20mm] (A) at (8.0,0.0) {$\mat{A}(\omega)$};
  \draw[->, thick, simInk] (sum.east) -- (A.west);
  \node[font=\scriptsize\itshape, align=center] at (8.0,-1.1)
    {one system operator\\\emph{per frequency}};
  \node[font=\scriptsize, text=simBlue, align=center] at (5.0,-2.3)
    {\code{assemble\_frequency\_operator}\\$=$ \code{linear\_combination}};
\end{tikzpicture}
\caption[The system operator as an operator-valued linear combination]{The
frequency-dependent system operator is an \emph{operator-valued linear
combination} (\cref{ch:combinators}) of the fixed basis. The three operators
$\Kop,\Cop,\Mop$ are assembled once and never change. The scalar weights
$1,\,i\omega,\,-\omega^2$ \emph{do} change with $\omega$. Summing
weight-times-operator yields the system operator $\mat{A}(\omega)$---a fresh
operator at every frequency, formed from a fixed basis. This is the same
\code{linear\_combination} combinator that sums tensors in BLAS-1; here its
operands are \emph{operators} and its scalars are \emph{complex}.}
\label{fig:opcombo}
\end{figure}

\subsection{The forming step is a linear combination}

Forming $\mat{A}(\omega)$ from the basis is not a new idea. It is the
\code{linear\_combination} combinator from \cref{ch:combinators}---the same one
that builds $\alpha\vx + \beta\vy$ in BLAS-1---specialized to the case where the
\emph{operands are operators} rather than tensors and the \emph{scalars are
complex} (\cref{fig:opcombo}). Spelled out as the framework verb
\code{assemble\_frequency\_operator}:
\begin{lstlisting}[style=l4style]
-- form the per-omega system operator from the fixed basis
assemble_frequency_operator :: FrequencyOperatorFamily -> Scalar -> LinOp
assemble_frequency_operator fam omega =
  linear_combination                  -- operator-operand corner; complex scalars
    [ (1,            fam.K)           --  K   with weight 1
    , (1i * omega,   fam.C)           --  C   with weight  i*omega
    , (-(omega^2),   fam.M)           --  M   with weight -omega^2
    , (1,            fam.A2 omega) ]  --  A2(omega) with weight 1
\end{lstlisting}
There is no separate fold algebra here---the body is literally
\code{linear\_combination} at term-list length four, with the operand monoid being
operator addition (commutative, identity the zero operator) instead of tensor
addition. The framework gives it its own name only because it is the named verb the
driven sweep reaches for; underneath, it is the BLAS-1 combinator the whole book
already trusts. The system operator $\mat{A}(\omega)$ is \emph{affine in $\omega$}
in the monomial part ($1, i\omega, -\omega^2$) modulo the extra term
$\mat{A}_2$; that affine structure is exactly what makes a reduced-order model
possible (\cref{sec:cost}), though the uniform sweep does not exploit it.

\begin{notationbox}
The system operator carries a complex element type because two of its three
weights are complex: $i\omega$ is imaginary and multiplies the damping $\Cop$,
$-\omega^2$ is real and multiplies the mass $\Mop$. So even when $\Kop,\Cop,\Mop$
are real operators, $\mat{A}(\omega)$ is a \emph{complex} operator, and the field
$\vE$ it solves for is a complex (phasor) field: its magnitude is the amplitude of
the oscillation and its argument is the phase. This is why the driven solve needs
GMRES over $\Complex^N$, not CG over $\Real^N$ (\cref{sec:solve}). The complex
arithmetic is not an implementation detail; it carries the phase information that
\emph{is} the answer.
\end{notationbox}

\section{Solve: the operator-varying map (the hoist, negated)}
\label{sec:solve}

Now the sweep. We are given a list of frequencies
$[\omega_0, \omega_1, \dots, \omega_{m-1}]$---the band, sampled---and at each one
we must solve \cref{eq:driven-system} for the field $\vE_k = \vE(\omega_k)$. The
solves are \emph{independent}: the field at $\omega_3$ does not depend on the
field at $\omega_2$, because each is the steady-state response to its own
frequency's excitation. Independence means the sweep is a \emph{map}, not a fold
(\cref{ch:situating}): the frequencies could be solved in any order, or on
separate machines, and the answers would be identical.

What makes it the operator-\emph{varying} map---corner~2 of \cref{ch:situating},
not corner~1---is that the operator on the left-hand side is not the same from one
iteration to the next. At each frequency we must:
\begin{enumerate}[nosep]
  \item \textbf{rebuild} the system operator $\mat{A}(\omega_k)$ from the fixed
  basis (\cref{sec:assemble});
  \item \textbf{re-capture} it into the solver---the call that tells the Krylov
  solver ``\emph{this} is the matrix you are about to invert'';
  \item \textbf{form} the right-hand side $\vb(\omega_k)$ from the port
  excitation at this frequency;
  \item \textbf{solve} $\mat{A}(\omega_k)\,\vE_k = \vb(\omega_k)$ by preconditioned
  GMRES (\cref{ch:gmres}).
\end{enumerate}
The framework names this combinator \code{frequency\_sweep}:
\begin{lstlisting}[style=l4style]
-- the operator-VARYING map: rebuild the operator INSIDE the loop
frequency_sweep :: FrequencyOperatorFamily -> [Scalar] -> [Tensor]
frequency_sweep fam omegas =
  map solve_at omegas               -- independent: a map, not a fold
  where
    solve_at omega =
      let a = assemble_frequency_operator fam omega   -- (1) rebuild A(omega)
          b = excitation fam omega                    -- (3) per-omega RHS
      in  ksp_solve a b                               -- (2,4) capture + GMRES solve
\end{lstlisting}

\subsection{The one line that is the whole difference}

Set the driven sweep beside electrostatics and almost everything is shared: both
assemble a basis once, both run a family of independent solves, both produce a
family of fields, both feed a rank-2 reduction. The single structural difference
is the \emph{position of the operator-capture call} relative to the loop
(\cref{fig:hoist}).

\begin{lstlisting}[style=l4style]
-- ELECTROSTATIC (corner 1): operator-capture HOISTED OUTSIDE the loop
solve_family k rhss =                    -- k assembled once, captured once
  let solver = ksp_setup k               -- <<-- SetOperators OUTSIDE
  in  map (\b -> ksp_run solver b) rhss   -- map: same operator every solve

-- DRIVEN (corner 2): operator-capture INSIDE the loop
frequency_sweep fam omegas =
  map (\omega ->
        let a      = assemble_frequency_operator fam omega
            solver = ksp_setup a          -- <<-- SetOperators INSIDE
        in  ksp_run solver (excitation fam omega))
      omegas                              -- map: fresh operator every solve
\end{lstlisting}

In electrostatics, \code{ksp\_setup}---the call that factors the preconditioner and
hands the operator to the solver---happens \emph{once}, before the map, because the
operator $\Kop$ is the same for every right-hand side. We \emph{hoist} it out of
the loop. In the driven sweep, \code{ksp\_setup} happens \emph{inside} the map,
once per frequency, because the operator $\mat{A}(\omega)$ is a function of the
loop index. We \emph{cannot} hoist it.

\begin{figure}[p]
\centering
\begin{tikzpicture}[font=\footnotesize, >=Stealth]
  % ============== LEFT: electrostatic, hoisted ==============
  \begin{scope}[shift={(0,0)}]
    \node[font=\small\bfseries, text=simBlue] at (1.5,4.3)
      {Electrostatic (corner 1)};
    \node[font=\scriptsize\itshape, text=simGray] at (1.5,3.9)
      {assemble once, \emph{hoist}};
    % the operator capture, outside
    \node[opbox, minimum width=30mm, fill=simBgBlue, draw=simBlue]
      (setup1) at (1.5,3.2) {\code{SetOperators}($\Kop$)};
    \node[font=\scriptsize, text=simBlue] at (1.5,2.65) {(once, OUTSIDE)};
    % the loop box
    \node[draw=simGray, densely dashed, rounded corners,
          minimum width=42mm, minimum height=24mm, fill=simBgGray]
      (loop1) at (1.5,0.6) {};
    \node[font=\scriptsize, text=simGray, anchor=north west]
      at (loop1.north west) {\ for each $\vb_i$:};
    \node[opbox, minimum width=26mm] (sv1) at (1.5,0.4)
      {\code{ksp\_run}($\vb_i$)};
    \node[font=\scriptsize\itshape] at (1.5,-0.4) {same $\Kop$ each pass};
    \draw[->, thick, simInk] (setup1.south) -- (loop1.north);
    % loop back arrow
    \draw[->, simGray, densely dashed]
      (loop1.east) ++(0,-0.4) .. controls +(0.9,0) and +(0.9,0) ..
      ($(loop1.east)+(0,0.6)$);
  \end{scope}

  % divider
  \draw[simGray, thin] (4.1,-1.0) -- (4.1,4.5);

  % ============== RIGHT: driven, inside ==============
  \begin{scope}[shift={(5.3,0)}]
    \node[font=\small\bfseries, text=simRust] at (1.7,4.3)
      {Driven (corner 2)};
    \node[font=\scriptsize\itshape, text=simGray] at (1.7,3.9)
      {basis once, rebuild \emph{inside}};
    % the basis, outside (only the basis is hoisted)
    \node[opbox, minimum width=34mm, fill=simBgBlue, draw=simBlue]
      (basis2) at (1.7,3.2) {assemble basis $\{\Kop,\Cop,\Mop\}$};
    \node[font=\scriptsize, text=simBlue] at (1.7,2.65) {(once, OUTSIDE)};
    % the loop box
    \node[draw=simRust, densely dashed, rounded corners,
          minimum width=46mm, minimum height=30mm, fill=simBgRust]
      (loop2) at (1.7,0.1) {};
    \node[font=\scriptsize, text=simRust, anchor=north west]
      at (loop2.north west) {\ for each $\omega_k$:};
    \node[opbox, minimum width=36mm, fill=white] (form2) at (1.7,0.95)
      {form $\mat{A}(\omega_k)=\Kop{+}i\omega_k\Cop{-}\omega_k^2\Mop$};
    \node[opbox, minimum width=36mm, fill=simBgRust, draw=simRust]
      (setup2) at (1.7,0.1) {\code{SetOperators}($\mat{A}(\omega_k)$)};
    \node[font=\scriptsize, text=simRust] at (1.7,-0.45) {(every pass, INSIDE)};
    \node[opbox, minimum width=36mm, fill=white] (sv2) at (1.7,-1.05)
      {\code{ksp\_run}($\vb(\omega_k)$)};
    \draw[->, simInk] (form2.south) -- (setup2.north);
    \draw[->, simInk] (setup2.south) -- (sv2.north);
    \draw[->, thick, simInk] (basis2.south) -- (loop2.north);
    \draw[->, simRust, densely dashed]
      (loop2.east) ++(0,-0.4) .. controls +(0.9,0) and +(0.9,0) ..
      ($(loop2.east)+(0,1.0)$);
  \end{scope}
\end{tikzpicture}
\caption[The hoist, and its negation]{The chapter's centerpiece: the operator
capture, hoisted versus not. \emph{Left}---electrostatics assembles $\Kop$ once
and captures it into the solver \emph{outside} the loop; every pass of the map
solves a new right-hand side against the \emph{same} operator. \emph{Right}---the
driven sweep assembles only the \emph{basis} $\{\Kop,\Cop,\Mop\}$ outside the
loop, then, \emph{inside} the loop, forms a fresh $\mat{A}(\omega_k)$ and captures
it (\code{SetOperators}) before each solve. The capture call moves from outside
the loop to inside it. That single move is the entire structural difference
between the two analyses---and moving it the wrong way does not slow the code
down, it computes the wrong answer.}
\label{fig:hoist}
\end{figure}

\begin{keyidea}
\textbf{Driven $=$ the anchor with the operator rebuilt \emph{inside} the sweep,
plus a port-projection reduction.} Electrostatics hoists the operator capture
\emph{outside} the loop (corner~1, the fixed-operator map). The driven analysis
moves that one call \emph{inside} the loop (corner~2, the operator-varying map),
because the system operator $\mat{A}(\omega) = \Kop + i\omega\Cop - \omega^2\Mop$
is a function of the loop index. The fixed \emph{basis} $\{\Kop,\Cop,\Mop\}$ is
still assembled once; only the \emph{combination} is rebuilt per frequency. Cap
the sweep with the S-parameter port-projection instead of the Gram, and you have
the entire driven driver. The position of one capture call relative to one loop is
load-bearing: hoist it and you solve every frequency against the operator of the
first one.
\end{keyidea}

\begin{pitfall}
A well-meaning optimizer sees the \code{SetOperators} call inside the loop, notices
it is ``expensive'' (it refactors the preconditioner), and hoists it out to ``save
work.'' In electrostatics that is correct and free. In the driven sweep it is a
\emph{silent wrong answer}: every frequency is now solved against
$\mat{A}(\omega_0)$, the operator of the first frequency, so the computed
$\vE_k$ is the steady-state response of the wrong physical system. Nothing crashes;
the numbers are simply wrong, and wrong in a way that looks plausible (a smooth
curve, just not the device's curve). The lesson of \cref{ch:situating} stated
operationally: the loop position of one setup call is part of the algorithm, not a
performance knob.
\end{pitfall}

\subsection{Why the solve is GMRES, not CG}

The driven system operator $\mat{A}(\omega) = \Kop + i\omega\Cop - \omega^2\Mop$ is
\emph{complex} and \emph{non-Hermitian}: the $i\omega\Cop$ term is purely
imaginary, and the $-\omega^2\Mop$ term subtracts mass from stiffness, so the
operator is indefinite as well. None of the structure CG relies on (\cref{ch:cg})
survives: there is no symmetric positive-definite form, hence no short
recurrence, hence no CG. This is exactly the regime GMRES was built for
(\cref{ch:gmres})---a general nonsymmetric, possibly complex non-Hermitian
operator---and the driven solve is its canonical application. Each frequency's
solve is a full preconditioned GMRES (or FGMRES, if the preconditioner itself
varies), and because the operator changes every frequency, the preconditioner must
be rebuilt or refreshed every frequency too (\cref{sec:cost}).

\section{The composition: \texttt{driven} end to end}

We can now write the whole driver as one composition, in the framework's
vocabulary. It is three stages: assemble the basis once, sweep with the operator
rebuilt inside, reduce to S-parameters (\cref{fig:pipeline}).
\begin{lstlisting}[style=l4style]
-- the whole driven driver: config in, frequency response out
driven :: DrivenConfig -> FrequencyResponse
driven cfg =
  let space  = nd_space cfg                                    -- H(curl) space
      fam    = { K  = fe_assemble space [ curl_curl (reluctivity cfg) ]  -- assemble
               , C  = fe_assemble space [ damping cfg ]                  --   basis
               , M  = fe_assemble space [ mass (permittivity cfg) ]      --   ONCE
               , A2 = \w -> extra_system_matrix space cfg w }
      omegas = sample_frequencies cfg                          -- the swept band
      es     = frequency_sweep fam omegas    -- operator-VARYING map: A(w) rebuilt inside
  in  sparameter_reduce (ports cfg) es       -- port-projection reduction -> S(w)
\end{lstlisting}
Read it as a pipeline composition:
\[
  \texttt{driven}
    \;=\;
  \texttt{sparameter\_reduce}
    \;\circ\;
  \texttt{frequency\_sweep}
    \;\circ\;
  \texttt{fe\_assemble}^{\times 3}.
\]
The first stage (\code{fe\_assemble}, three times) is the once-only basis; the
middle stage (\code{frequency\_sweep}) is the operator-varying solve map; the last
stage (\code{sparameter\_reduce}) is the projection readout we turn to next. Each
stage is firm, and each links down to its own chapter: the assemble fold to
\cref{ch:assembly}, the combination to \cref{ch:combinators}, the solve to
\cref{ch:gmres}, the reduction to \cref{ch:reductions}.

\begin{figure}[t]
\centering
\begin{tikzpicture}[node distance=7mm and 9mm, font=\footnotesize, >=Stealth]
  \node[data] (cfg) {config};
  \node[stage, right=of cfg, align=center] (asm)
    {assemble basis\\$\{\Kop,\Cop,\Mop\}$\\\scriptsize(once)};
  \node[stage, right=of asm, fill=simBgRust, draw=simRust, align=center] (sweep)
    {\textbf{frequency\_sweep}\\\scriptsize per $\omega$:};
  \node[stage, right=of sweep, align=center] (red)
    {sparameter\_\\reduce};
  \node[product, right=of red] (S) {$S(\omega)$};
  \draw[flow] (cfg) -- (asm);
  \draw[flow] (asm) -- node[above,font=\scriptsize]{\code{fam}} (sweep);
  \draw[flow] (sweep) -- node[above,font=\scriptsize]{$[\vE_k]$} (red);
  \draw[flow] (red) -- (S);
  % the per-omega inner body, drawn below the sweep box
  \node[opbox, below=10mm of sweep, fill=white, align=center] (inner)
    {form $\mat{A}(\omega)$ $\to$ \code{SetOperators} $\to$ GMRES};
  \draw[refedge] (sweep.south) -- (inner.north);
  \node[font=\scriptsize\itshape, text=simRust, below=1.5mm of inner]
    {operator rebuilt + re-captured \emph{inside}};
\end{tikzpicture}
\caption[The driven pipeline]{The driven pipeline. The basis
$\{\Kop,\Cop,\Mop\}$ is assembled once. \code{frequency\_sweep} maps over the
swept band; its per-$\omega$ body (drawn below) forms the system operator
$\mat{A}(\omega)$, re-captures it with \code{SetOperators}, and solves by GMRES.
The solved fields $[\vE_k]$ feed \code{sparameter\_reduce}, the port-projection
that yields the scattering matrix $S(\omega)$. The red box marks the one stage
whose operator changes each step---the entire difference from electrostatics.}
\label{fig:pipeline}
\end{figure}

\begin{remark}[Uniform sweep versus adaptive ROM]
The composition above is the \emph{uniform} sweep: a fixed list of frequencies,
each solved fully and independently---an honest map. Palace also offers an
\emph{adaptive} sweep, which builds a projection-based reduced-order model (a PROM)
and greedily adds the next sample frequency where the model's error is largest. The
adaptive sweep is \emph{not} a map: each new frequency choice depends on the model
built from all previous solves, so it is a state-generated \emph{fold}
(\cref{ch:situating}, \cref{ch:transient})---the running model is threaded through
the loop. We mention it only to place it: the uniform sweep is the composed shape
this chapter studies, and it is the operator-varying map. The affine-in-$\omega$
structure of $\mat{A}(\omega)$ (\cref{sec:assemble}) is exactly what makes the
reduced-order model cheap to evaluate, but the uniform sweep does not exploit it.
\end{remark}

\section{Reduce: the port-projection}
\label{sec:reduce}

The solve gives us a family of complex fields, one per swept frequency (and, when
several ports are driven, one column per driven port). The reduction collapses
each solved field into the scattering parameters the user asked for. This is the
\emph{rank-2 port-projection}---reduction shape~2 of \cref{ch:situating} and
\cref{ch:reductions}---and it is emphatically \emph{not} a Gram.

The port modes $\vs_i$ themselves do not come from the driven solve. They come
from the \emph{boundary-mode} analysis (\cref{ch:waveports}), which solves a small
two-dimensional eigenproblem on each port's cross-section to find the field
pattern that propagates there. The driven reduction \emph{consumes} those modes as
the basis it projects onto. Given the solved field $\vE_j$ for drive port $j$ and
the receiver modes $\vs_i$, each scattering entry is
\begin{equation}
  \boxed{\;S_{ij}
    \;=\;
  \sigma_{ij}\,\bigl(\,\inner{\vs_i}{\vE_j(\omega)} \;-\; \delta_{ij}\,\bigr)\;}
  \label{eq:sparam}
\end{equation}
where each piece earns its place:
\begin{itemize}[nosep]
  \item $\inner{\vs_i}{\vE_j(\omega)}$ is the \emph{linear projection} of the
  drive-$j$ field onto receiver mode $i$---how much of the solved field looks like
  the wave that propagates out of port $i$. For a lumped port this is the inner
  product $\vs_i\cdot\vE_j$; for a wave port it is the modal overlap
  $(\vE_j\times\vH_i^\star)\cdot\vn$ integrated over the port face. It is
  \emph{linear} in the field $\vE_j$.
  \item $\delta_{ij}$ is the Kronecker self-term: $1$ when $i=j$, $0$ otherwise.
  On the diagonal it subtracts the \emph{incident} wave from the total field,
  leaving only the \emph{reflected} part. Without it, $S_{11}$ would report the
  whole field at port~1 (incident plus reflected); with it, $S_{11}$ reports the
  reflection alone, which is what ``return loss'' means.
  \item $\sigma_{ij}$ is the port-kind \emph{scale}: an impedance ratio
  $\sqrt{R_j/R_i}$ for lumped ports (generalized scattering parameters, which
  normalize for unequal port impedances), or a de-embedding phase
  $e^{ik_{n,i}d_i}\,e^{ik_{n,j}d_j}$ for wave ports (which shift the reference
  plane from the mesh boundary to the physical port location).
\end{itemize}

\begin{figure}[t]
\centering
\begin{tikzpicture}[font=\footnotesize, >=Stealth]
  % field-space axes
  \draw[->, simGray] (0,0) -- (0,3.4);
  \draw[->, simGray] (0,0) -- (5.0,0);
  % port mode basis vectors
  \draw[->, very thick, simGreen] (0,0) -- (3.6,0)
    node[right, font=\scriptsize, text=simGreen]{$\vs_1$ (port 1 mode)};
  \draw[->, very thick, simGreen] (0,0) -- (0,2.8)
    node[above, font=\scriptsize, text=simGreen]{$\vs_2$ (port 2 mode)};
  % solved field
  \draw[->, very thick, simRust] (0,0) -- (2.6,2.0)
    node[midway, above left, font=\scriptsize, text=simRust]{$\vE_1(\omega)$};
  % projections (dashed)
  \draw[densely dashed, simBlue] (2.6,2.0) -- (2.6,0)
    node[midway, right, font=\scriptsize, text=simBlue]{};
  \draw[densely dashed, simBlue] (2.6,2.0) -- (0,2.0);
  \node[font=\scriptsize, text=simBlue] at (2.6,-0.32)
    {$\inner{\vs_1}{\vE_1}$};
  \node[font=\scriptsize, text=simBlue, anchor=east] at (-0.1,2.0)
    {$\inner{\vs_2}{\vE_1}$};
  % the entries
  \node[font=\footnotesize, align=left, text=simInk] at (6.9,2.2)
    {$S_{11}=\sigma_{11}(\inner{\vs_1}{\vE_1}-1)$};
  \node[font=\scriptsize, text=simGray] at (6.9,1.75)
    {(reflection: self-term subtracted)};
  \node[font=\footnotesize, align=left, text=simInk] at (6.9,1.0)
    {$S_{21}=\sigma_{21}\,\inner{\vs_2}{\vE_1}$};
  \node[font=\scriptsize, text=simGray] at (6.9,0.55)
    {(transmission: no self-term)};
\end{tikzpicture}
\caption[The port projection]{The port-projection reduction. The solved field
$\vE_1(\omega)$ (driving port~1) lives in field space; the port modes $\vs_1,\vs_2$
(from the boundary-mode analysis, \cref{ch:waveports}) are the directions we
project onto. Projecting $\vE_1$ onto $\vs_1$ and subtracting the incident
self-term gives the reflection $S_{11}$; projecting onto $\vs_2$ (no self-term,
since $2\neq1$) gives the transmission $S_{21}$. The projection is \emph{linear}
in the field and the two index families ($\vs_i$ versus $\vE_j$) play different
roles---this is why it is a projection, not a Gram.}
\label{fig:proj}
\end{figure}

In the framework this reduction is the verb \code{sparameter\_reduce}, applied once
per swept frequency:
\begin{lstlisting}[style=l4style]
-- project each solved field onto each receiver port mode -> the S-matrix
sparameter_reduce :: [PortMode] -> [(Int, Tensor)] -> Matrix
sparameter_reduce ports family =
  matrix_from_columns
    [ [ scale ports i j * (project (ports !! i) e - selfterm i j)  -- entry S_ij
        | i <- [0 .. p-1] ]            -- over receiver ports i
      | (j, e) <- family ]            -- one column per drive port j
  where
    p              = length ports
    project s e    = port_dot s e               -- linear functional  s_i . E
    selfterm i j   = if i == j then 1 else 0    -- the -1 self-reflection (diagonal)
    scale ports i j = port_scale (ports!!i) (ports!!j)  -- impedance / de-embed
\end{lstlisting}

\subsection{Why it is rank-2 like a Gram but deliberately not a Gram}

The S-matrix and the capacitance matrix (\cref{ch:reductions}) both come out as
$p\times p$ arrays indexed by a port-like family. That shared \emph{output rank}
is the whole of their resemblance, and \cref{ch:situating} already warned that
reading more into it is the single most over-unified mistake the design could
make. The driven reduction makes the distinction sharp.

\begin{pitfall}
\textbf{A matrix-shaped result is not a Gram.} The Gram reduction
(capacitance, inductance) pairs a family \emph{with itself} through an operator,
$\vx_j^{\!\top}\Kop\vx_i$---a \emph{bilinear} form, symmetric by construction, with
the verb computing only the upper triangle. The S-parameter reduction projects one
family (the solved fields $\vE_j$) onto a \emph{different} basis (the port modes
$\vs_i$)---a \emph{linear} functional, two families in different roles, every entry
computed independently. Three facts forbid merging them: (1) the Gram is a
self-pairing, the projection is not; (2) the Gram is symmetric by construction,
the projection's near-symmetry $S_{ij}\approx S_{ji}$ is reciprocity \emph{physics}
when the medium is reciprocal, never a construction the reduction imposes---there
is no \code{symmetric\_from\_upper} in its body; (3) the Gram is real and
energy-valued, the projection is complex and carries the inhomogeneous
$-\delta_{ij}$ self-term on its diagonal. Forcing one verb would smuggle a false
symmetry into the S-matrix or burden the Gram with machinery it does not need.
\emph{Sharing a shape is not sharing a meaning.}
\end{pitfall}

\section{Reading the answer: return loss and insertion loss}

The product is the scattering matrix as a function of frequency, $S(\omega)$, and
microwave engineers read it through two standard plots, both in decibels
(\cref{fig:splot}). The \emph{return loss} is $\abs{S_{11}}$ in dB---how much power
reflects off port~1; near a perfect match it dives toward $-\infty$ (a sharp dip),
because almost nothing reflects. The \emph{insertion loss} is $\abs{S_{21}}$ in
dB---how much power transmits from port~1 to port~2; in a passband it sits near
$0\,$dB (almost everything gets through), and in a stopband it plunges (the device
blocks the signal). A bandpass filter shows $\abs{S_{21}}$ flat and high across its
passband and low outside it, with $\abs{S_{11}}$ dipping sharply wherever the match
is good. The whole point of the simulation is to predict these curves before the
device is fabricated.

\begin{figure}[t]
\centering
\begin{tikzpicture}
\begin{axis}[
  width=0.84\linewidth, height=6.0cm,
  xlabel={frequency $f$ (GHz)}, ylabel={magnitude (dB)},
  xmin=3, xmax=7, ymin=-45, ymax=3,
  ytick={0,-10,-20,-30,-40},
  xtick={3,4,5,6,7},
  legend pos=south east, legend cell align=left,
  grid=both, grid style={simGray!20},
  tick label style={font=\footnotesize}, label style={font=\footnotesize},
  legend style={font=\footnotesize},
]
% |S21| insertion loss: passband near 0 dB around 5 GHz, rolling off outside
\addplot[very thick, simGreen, smooth, samples=120, domain=3:7]
  { -2 - 30*( (x-5)/1.0 )^2 / (1 + 0.55*((x-5))^2 ) - 0.5*(x-5)^2 };
\addlegendentry{$|S_{21}|$ (insertion loss)}
% |S11| return loss: deep dip near the match frequency 5 GHz
\addplot[very thick, simRust, smooth, samples=160, domain=3:7]
  { -3 - 32*exp( -((x-5)^2)/0.06 ) + 2.0*((x-5))^2 };
\addlegendentry{$|S_{11}|$ (return loss)}
% passband band marker
\draw[simBlue!50, densely dashed] (axis cs:4.4,3) -- (axis cs:4.4,-45);
\draw[simBlue!50, densely dashed] (axis cs:5.6,3) -- (axis cs:5.6,-45);
\node[font=\scriptsize, text=simBlue] at (axis cs:5.0,-43) {passband};
\end{axis}
\end{tikzpicture}
\caption[A swept S-parameter plot]{The product: a swept S-parameter plot for a
bandpass-like two-port. The insertion loss $\abs{S_{21}}$ (green) sits near
$0\,$dB across the passband---signal passes through---and rolls off on either
side. The return loss $\abs{S_{11}}$ (rust) dips sharply at the match frequency
near $5\,$GHz, where almost no power reflects (a good match), and rises away from
it. Each point on each curve is one converged GMRES solve at one frequency,
projected onto the port modes. Reading these two curves is reading the device.}
\label{fig:splot}
\end{figure}

\section{Cost: many independent solves, preconditioner reuse}
\label{sec:cost}

Because the sweep is a \emph{map}, the frequencies are embarrassingly parallel: a
hundred-frequency sweep is a hundred independent GMRES solves, and they can run on a
hundred machines with no communication. That is the good news, and it follows
directly from corner~2's independence.

The bad news is that each of those solves is a \emph{full} GMRES on a complex
indefinite operator, and---unlike electrostatics---we cannot amortize a single
factorization across the family, because the operator $\mat{A}(\omega)$ is
different every frequency. The preconditioner (\cref{ch:precond}) is what makes
each solve converge in a reasonable number of iterations, and rebuilding it from
scratch at every frequency is expensive. Two structural facts soften this:
\begin{itemize}[nosep]
  \item \textbf{Nearby frequencies have nearby operators.} Because
  $\mat{A}(\omega)$ depends smoothly on $\omega$, the operator at $\omega_{k+1}$ is
  close to the operator at $\omega_k$. A preconditioner built for $\omega_k$ is
  often a perfectly good preconditioner for $\omega_{k+1}$, so one can \emph{reuse}
  a preconditioner across a band of nearby frequencies and only rebuild it
  occasionally---trading a little extra GMRES iteration for a lot of saved
  factorization. When the preconditioner is allowed to drift this way, the inner
  solve becomes FGMRES (\cref{ch:gmres}), the flexible variant that tolerates a
  changing preconditioner.
  \item \textbf{The affine structure enables a reduced-order model.} The
  affine-in-$\omega$ form $\mat{A}(\omega) = \Kop + i\omega\Cop - \omega^2\Mop$
  means the solution $\vE(\omega)$ is a rational function of $\omega$, and a small
  number of full solves can fit a reduced-order model that predicts $\vE$ at the
  remaining frequencies cheaply. This is the adaptive sweep's payoff (the
  state-generated fold of the earlier remark); the uniform sweep leaves it on the
  table in exchange for simplicity and perfect parallelism.
\end{itemize}

\section{Two worked examples}

\begin{workedexample}{The operator-varying map, made concrete}{opvary}
We make the central claim---the operator changes every frequency---numerical, on a
deliberately tiny basis so the arithmetic is visible. Take the $2\times2$ real
basis operators
\[
  \Kop = \begin{pmatrix} 4 & -1 \\ -1 & 3 \end{pmatrix},\quad
  \Cop = \begin{pmatrix} 2 & 0 \\ 0 & 1 \end{pmatrix},\quad
  \Mop = \begin{pmatrix} 1 & 0 \\ 0 & 1 \end{pmatrix},
\]
assembled \emph{once}, and take $\mat{A}_2 \equiv 0$ for clarity. The system
operator is $\mat{A}(\omega) = \Kop + i\omega\Cop - \omega^2\Mop$. We form it at
three frequencies.

\medskip
\emph{At $\omega = 0$} (the static limit): the weights are $1,\,0,\,0$, so
\[
  \mat{A}(0) = \Kop = \begin{pmatrix} 4 & -1 \\ -1 & 3 \end{pmatrix}.
\]
At zero frequency the driven operator collapses to the electrostatic stiffness:
no damping, no mass. This is the anchor reappearing as a special case.

\medskip
\emph{At $\omega = 1$:} the weights are $1,\,i,\,-1$, so
\[
  \mat{A}(1)
    = \Kop + i\Cop - \Mop
    = \begin{pmatrix} 4-1 & -1 \\ -1 & 3-1 \end{pmatrix}
      + i\begin{pmatrix} 2 & 0 \\ 0 & 1 \end{pmatrix}
    = \begin{pmatrix} 3+2i & -1 \\ -1 & 2+i \end{pmatrix}.
\]

\emph{At $\omega = 2$:} the weights are $1,\,2i,\,-4$, so
\[
  \mat{A}(2)
    = \Kop + 2i\Cop - 4\Mop
    = \begin{pmatrix} 4-4 & -1 \\ -1 & 3-4 \end{pmatrix}
      + i\begin{pmatrix} 4 & 0 \\ 0 & 2 \end{pmatrix}
    = \begin{pmatrix} 0+4i & -1 \\ -1 & -1+2i \end{pmatrix}.
\]

Three frequencies, \emph{three different operators}---and all three complex, all
three formed from the same fixed $\{\Kop,\Cop,\Mop\}$ by changing only the scalar
weights. Contrast electrostatics, where the operator is $\Kop$ at \emph{every}
solve. This is corner~1 versus corner~2 in two lines of arithmetic.

\medskip
\emph{One solve.} Take the excitation $\vb = (1,\,0)^{\!\top}$ and solve
$\mat{A}(1)\,\vE = \vb$. With
$\det\mat{A}(1) = (3+2i)(2+i) - (-1)(-1) = (4 + 7i) - 1 = 3 + 7i$,
\[
  \vE
    = \frac{1}{3+7i}
      \begin{pmatrix} 2+i & 1 \\ 1 & 3+2i \end{pmatrix}
      \begin{pmatrix} 1 \\ 0 \end{pmatrix}
    = \frac{1}{3+7i}\begin{pmatrix} 2+i \\ 1 \end{pmatrix}.
\]
Multiplying numerator and denominator by the conjugate $3-7i$ (with
$\abs{3+7i}^2 = 58$),
\[
  \vE
    = \frac{1}{58}\begin{pmatrix} (2+i)(3-7i) \\ (3-7i) \end{pmatrix}
    = \frac{1}{58}\begin{pmatrix} 13 - 11i \\ 3 - 7i \end{pmatrix}
    \approx \begin{pmatrix} 0.224 - 0.190\,i \\ 0.052 - 0.121\,i \end{pmatrix}.
\]
The solved field is genuinely complex: it has magnitude (amplitude) and phase, and
the phase is the steady-state lag of the field behind the drive---information that
a real-valued static solve simply does not carry. In the real driver this $2\times2$
direct solve is replaced by preconditioned GMRES on a million-dimensional complex
operator, but the structure is exactly this: rebuild $\mat{A}(\omega)$, then solve.
\end{workedexample}

\begin{workedexample}{A two-port S-parameter readout}{sparam}
Now the reduction. Suppose we have driven port~1 of a two-port device at one
frequency and solved for the field $\vE_1$. The boundary-mode analysis
(\cref{ch:waveports}) gave us the two port modes $\vs_1,\vs_2$, and the projections
of the solved field onto them came out as
\[
  \inner{\vs_1}{\vE_1} = 0.20 + 0.10\,i,
  \qquad
  \inner{\vs_2}{\vE_1} = 0.70 - 0.30\,i.
\]
Take lumped ports of equal impedance, so the scale is trivial,
$\sigma_{ij} = \sqrt{R_j/R_i} = 1$. Apply \cref{eq:sparam}.

\medskip
\emph{Reflection $S_{11}$} (receiver $1$, drive $1$, so $i=j$---subtract the
self-term):
\[
  S_{11}
    = \sigma_{11}\bigl(\inner{\vs_1}{\vE_1} - \delta_{11}\bigr)
    = (0.20 + 0.10\,i) - 1
    = -0.80 + 0.10\,i.
\]
\emph{Transmission $S_{21}$} (receiver $2$, drive $1$, so $i\neq j$---no
self-term):
\[
  S_{21}
    = \sigma_{21}\bigl(\inner{\vs_2}{\vE_1} - \delta_{21}\bigr)
    = (0.70 - 0.30\,i) - 0
    = 0.70 - 0.30\,i.
\]

\medskip
\emph{Read them.} The magnitudes are
\[
  \abs{S_{11}} = \sqrt{0.80^2 + 0.10^2} \approx 0.806,
  \qquad
  \abs{S_{21}} = \sqrt{0.70^2 + 0.30^2} \approx 0.762.
\]
In decibels, the return loss is
$20\log_{10}\abs{S_{11}} \approx -1.87\,$dB and the insertion loss is
$20\log_{10}\abs{S_{21}} \approx -2.36\,$dB. A return loss this shallow means a
\emph{poor} match---about $65\%$ of the incident power
($\abs{S_{11}}^2 \approx 0.65$) reflects straight back---so this frequency is not
in a passband. Notice the load-bearing role of the self-term: without the
$-\delta_{11}$, $S_{11}$ would have been the raw projection $0.20+0.10\,i$, with
magnitude $\approx 0.22$, reporting a \emph{good} match (low reflection)---the
exact opposite conclusion. The $-1$ is what converts ``total field at the driven
port'' into ``reflected field at the driven port,'' and it is the difference
between a right answer and a plausible wrong one.

\medskip
Sweeping this readout across the band---one solved field, one projection, per
frequency---traces the $\abs{S_{11}}$ and $\abs{S_{21}}$ curves of
\cref{fig:splot}. That sweep is the whole driven driver.
\end{workedexample}

\summarysection
The driven analysis computes the frequency response of a microwave device: feed a
time-harmonic signal into the ports, sweep the frequency band, and read off the
scattering matrix $S(\omega)$. It is the anchor (\cref{ch:electrostatics}) with one
corner swapped. \emph{Assemble} builds a fixed basis $\{\Kop,\Cop,\Mop\}$ once, but
the system operator $\mat{A}(\omega) = \Kop + i\omega\Cop - \omega^2\Mop$
$[+\mat{A}_2(\omega)]$ is a frequency-dependent linear combination of that basis
(\code{assemble\_frequency\_operator}, a specialization of \code{linear\_combination}
with operator operands and complex scalars). \emph{Solve} is the operator-varying
map \code{frequency\_sweep}: at each frequency, rebuild $\mat{A}(\omega)$, re-capture
it into the solver \emph{inside} the loop (the exact negation of electrostatics'
hoist), and solve the complex non-Hermitian system by preconditioned GMRES
(\cref{ch:gmres}). The solves are independent (a map, hence parallel), but each is a
full GMRES, so preconditioner reuse across nearby frequencies matters. \emph{Reduce}
is the rank-2 port-projection \code{sparameter\_reduce}:
$S_{ij} = \sigma_{ij}(\inner{\vs_i}{\vE_j} - \delta_{ij})$---project each solved
field onto each port mode (from the boundary-mode analysis, \cref{ch:waveports}),
subtract the self-reflection on the diagonal, scale for impedance or de-embedding.
It produces a matrix like the Gram but is deliberately \emph{not} a Gram: a linear
projection of two distinct families, complex, asymmetric. The whole driver is
$\texttt{driven} = \texttt{sparameter\_reduce}\circ\texttt{frequency\_sweep}\circ
\texttt{fe\_assemble}^{\times3}$, and reading $\abs{S_{11}}$ (return loss) and
$\abs{S_{21}}$ (insertion loss) across the band is reading the device.

\furtherreading
The scattering-parameter formalism, ports, impedance normalization, and the
network-analyzer measurement this chapter simulates are the foundation of
microwave engineering; Pozar's text~\citep{pozar2011} is the standard treatment and
the place to deepen the physical reading of $S(\omega)$, return loss, insertion
loss, and de-embedding. The finite-element formulation of the time-harmonic
curl--curl system on $\Hcurl$, including wave-port boundary conditions and the modal
projection that defines the port modes $\vs_i$, is developed at length by
Jin~\citep{jin2014}. The complex, non-Hermitian, indefinite linear solve at each
frequency---and the GMRES/FGMRES machinery that handles it, together with
preconditioner reuse across a parameter sweep---is the subject of Saad's account of
iterative methods~\citep{saad2003}, on which \cref{ch:gmres} and \cref{ch:precond}
both draw. The boundary-mode analysis that produces the port modes is
\cref{ch:waveports}; the state-generated fold that the adaptive sweep becomes is the
shape of \cref{ch:transient}.

\exercisessection
\begin{enumerate}[nosep]
  \item \textbf{The one moved line.} In one sentence, name the single operation
  whose position---inside the frequency loop versus outside it---is the entire
  structural difference between the driven sweep and electrostatics. Explain why
  hoisting it out of the driven loop changes the \emph{answer}, not merely the
  speed.
  \item \textbf{Form the operator.} Using the basis of
  \cref{wex:opvary} ($\Kop,\Cop,\Mop$ as given, $\mat{A}_2\equiv0$), write down
  $\mat{A}(\omega)$ at $\omega = \tfrac12$ and at $\omega = 3$. Confirm that the two
  operators differ and that both are complex. Which weight is responsible for the
  imaginary part, and which for making the operator indefinite at large $\omega$?
  \item \textbf{Solve at a frequency.} With $\mat{A}(1)$ from \cref{wex:opvary} and
  the excitation $\vb = (0,\,1)^{\!\top}$, solve $\mat{A}(1)\vE = \vb$ by hand
  (the $2\times2$ inverse is given in the example). Report the magnitude and phase
  of each component of $\vE$.
  \item \textbf{The self-term.} In \cref{wex:sparam}, recompute $S_{11}$ and its
  magnitude \emph{without} the $-\delta_{11}$ self-term. Show that omitting it flips
  the physical conclusion from ``poor match'' to ``good match,'' and explain in one
  sentence what the $-1$ physically subtracts.
  \item \textbf{Projection, not Gram.} Give two structural reasons the S-parameter
  matrix is not a Gram matrix (\cref{ch:reductions}). For each reason, name the
  feature of \code{sparameter\_reduce}'s body that a Gram reduction would lack or
  forbid.
  \item \textbf{Why GMRES.} The driven system operator
  $\mat{A}(\omega) = \Kop + i\omega\Cop - \omega^2\Mop$ is solved by GMRES, not CG.
  Identify which two terms break the symmetric-positive-definite structure CG
  requires, and say what each one does to the operator (one makes it complex
  non-Hermitian; one makes it indefinite).
  \item \textbf{Parallelism and cost.} The frequency sweep is a map, so its solves
  are independent and parallel---yet the chapter calls each one expensive. Explain
  the tension: what can the driven sweep \emph{not} amortize across the family that
  electrostatics \emph{can}, and name the two mitigations the chapter offers.
  \item \textbf{Place it on the grid.} Using the driver $\times$ framework grid of
  \cref{ch:situating}, describe the driven analysis as ``electrostatics with these
  cells changed.'' Exactly which stage-cells (function space, assemble, solve
  corner, reduction) differ, and which stay the same?
\end{enumerate}
